\documentclass[10pt,a4paper]{amsart}
\usepackage[margin=19mm]{geometry}
\usepackage{amsmath,amssymb,mathtools}
\usepackage[scr=rsfs,cal=euler]{mathalfa}
\usepackage{xfrac}
\usepackage[unicode,hidelinks,bookmarks=false]{hyperref}
\newcommand{\ie}{i.e.\ }
\newcommand{\cf}{cf.\ }
\newcommand{\resp}{respectively\ }
\newcommand{\Subsection}{\subsection}
\providecommand{\nobreakdash}{\nobreak\hbox{-}}
\setlength{\emergencystretch}{2em}
\multlinegap0pt
\usepackage{amssymb}
\usepackage{tikz}
\newtheorem{prop}{Proposition}
\newcommand{\E}{\mathsf{E}}
\newcommand{\TT}{\mathbb{T}}
\newcommand{\credal}{\mathscr{Q}}
\newcommand{\eval}{\mathfrak{e}}
\newcommand{\prior}{\mathsf{Q}}
\newcommand{\prob}{\mathsf{P}} 
\newcommand{\uprior}{\overline{\mathsf{Q}}}
\newcommand{\uuprob}{\mathbf{\overline{P}}}
\title{Regularized e-processes: anytime valid inference with knowledge-based efficiency gains}
\author{Ryan Martin}
\date{Source: arXiv:2410.01427v4 (CC BY 4.0)}
\begin{document}
\maketitle

\noindent\emph{Source and licence.} Regularized e-processes: anytime valid inference with knowledge-based efficiency gains. Version arXiv:2410.01427v4 (CC BY 4.0), distributed by arXiv under the Creative Commons Attribution 4.0 International licence, http://creativecommons.org/licenses/by/4.0/, as recorded in the arXiv licence metadata for this exact version and shown on the abstract page https://arxiv.org/abs/2410.01427v4. Full licence notice: \texttt{licenses/arxiv-2410.01427-CC-BY-4.0.txt}.\par\smallskip

\noindent\textit{Editorial scope.} Counted unit: the article's prop prop:dominates (source0.tex lines 1451--1454). The article's complete original proof of it is delivered with it (source0.tex lines 1456--1494, one \texttt{proof} environment of 39 lines). No mathematics is written, repaired or re-derived: the statement and the proof are the article's own lines, in the article's own notation, and no retained line is substituted or re-flowed.  Retained uncounted as the source-local context the proof needs: lines 423--426; lines 1417--1420; lines 1424--1427; lines 1439--1442.  Nothing was omitted from any retained span, so the package carries no omission record.  No retained line contains a citation, so the wrapper carries no bibliography and no bibliography entry.  No retained line contains a figure environment, an image-inclusion command or drawing code, so no image is delivered and the package declares no asset dependency.  Editorial formatting only, in addition: an AMS \texttt{amsart} preamble that restates the article's own declarations and macros used by the retained lines, namely the environments \texttt{prop} on separate counters, together with the standard packages those lines need.  The retained environments print in this document as Proposition 1 in order of appearance, whereas the article prints the same items with the numbering of its own full document.  The complete native source of the article as distributed in the arXiv e-print for this exact version is bundled unchanged as \texttt{sources/166149/source0.tex}.
\begin{equation}
\label{eq:reg.eprocess}
\eval^\text{reg}(\cdot, \theta) = \rho(\theta) \times \eval_\theta(\cdot), \quad \theta \in \TT.
\end{equation}

\begin{equation}
\label{eq:sovereign}
e \geq 1 \implies m(r, e) \geq r, \quad \text{for all $r \in (0,\infty)$}, 
\end{equation}

\begin{equation}
\label{eq:merging}
\uuprob\bigl[ m\{ \rho(\Theta), \eval_{\Theta}(Z) \} \bigr] \leq 1 \quad \text{for all e-processes $\eval$}. 
\end{equation}

\begin{align}
\uuprob \bigl[ m\{ \rho(\Theta), \eval_\Theta(Z) \} > m'\{ \rho(\Theta), \eval_\Theta(Z) \}, \; & \notag \\
\rho(\Theta) \geq 1, \, \eval_\Theta(Z) \geq 1 & \bigr] > 0, \quad \text{for all e-processes $\eval$}. \label{eq:uuprob.dominance}
\end{align}

\begin{prop}
\label{prop:dominates}
Given $\uprior$, fix a non-trivial, admissible regularizer $\rho$.  Then the product rule in Equation~\eqref{eq:reg.eprocess} isn't $\uuprob$-strictly weakly dominated by any other re-merging function. 
\end{prop}

\begin{proof}
The proof is by contradiction; that is, I'll assume that the product rule defined in Equation~\eqref{eq:reg.eprocess} is $\uuprob$-strictly weakly dominated in the sense above and show that this leads to a contradiction.  Let $m$ denote this assumed-to-exist dominant re-merging function. 

The assumed non-triviality of the regularizer $\rho$ implies existence of points $\theta$ such that $\rho(\theta) > 1$ and such that $\rho(\theta) \leq 1$, and neither of these sets have $\uprior$-probability 0.  Some of the $\theta$'s in the first set could have $\rho(\theta) = \infty$, but, those must have  $\uprior$-probability 0 for, otherwise, the property that $\uprior\,\rho \leq 1$ would be violated; recall that the assumed admissibility of $\rho$ means that $\uprior\,\rho=1$.  

From the assumed strict weak dominance of $m$, and from \eqref{eq:sovereign}, I can deduce the following bound for generic inputs $(r,e)$:
\begin{align*}
m(r,e) & = m(r,e) \, \bigl( 1_{r < 1, e < 1} + 1_{r \geq 1, e < 1} + 1_{r < 1, e \geq 1} + 1_{r \geq 1, e \geq 1} \bigr) \\
& \geq  m(r,e) \, 1_{r < 1, e < 1} + m(r,e) \, 1_{r \geq 1, e < 1} + r \, 1_{r < 1, e \geq 1} + re \, 1_{r \geq 1, e \geq 1},
\end{align*}
where the ``$r$'' factor in the third term is by \eqref{eq:sovereign} and the ``$re$'' factor in the fourth term is by the assumed weak dominance.  The first two terms depend on details of the particular choice of $m$ on the respective ranges of $(r,e)$, details that can't be controlled with only the information provided.  I can apply the trivial non-negativity bound, however, and, from the above display, conclude that 
\[ m(r,e) \geq r \, 1_{r < 1, e \geq 1} + re \, 1_{r \geq 1, e \geq 1}. \]
(In fact, with the input e-process to be constructed next, those two terms I lower-bounded by 0 would typically be equal to 0, so the above inequality isn't loose.)  The not-strict equality in the above display is pointwise in $(r,e)$, i.e., I can't rule out equality above for any given pair $(r,e)$.  But a pointwise lower bound isn't the goal---I'm aiming for a lower bound in upper expectation.  For this latter goal, I'll apply \eqref{eq:uuprob.dominance} to flip the not-strict inequality ``$\geq$'' in the above display a strict inequality; more on this below.  

Towards establishing a contradiction, I only need to produce one example of an input e-process such that the conclusion is problematic, and I'll do this with an incredibly simple e-process.  Specifically, I define the input e-process $\eval^\star$ as follows:
\begin{itemize}
\item if $\theta$ is such that $\rho(\theta) < 1$, then $\eval_\theta^\star(Z) \equiv 1$, and 
\vspace{-2mm}
\item if $\theta$ is such that $\rho(\theta) \geq 1$, then  
\[ \eval_\theta^\star(Z) = \begin{cases} 2 & \text{if $Z \in \mathcal{E}_\theta$} \\ 0 & \text{otherwise}, \end{cases} \]
where $\mathcal{E}_\theta$ is an event with $\prob_\theta(\mathcal{E}_\theta) = \frac12$.
\end{itemize} 
It's easy to check that $\E_\theta\{ \eval_\theta^\star(Z) \} = 1$ for all $\theta$, so $\eval^\star$ is a genuine e-process.  What's important about this particular e-process, as it pertains to the present proof, is that 
\[ 1_{\rho(\theta) < 1, \eval_\theta^\star \geq 1} = 1_{\rho(\theta) < 1} \quad \text{and} \quad \eval_\theta^\star \, 1_{\rho(\theta) \geq 1, \eval_\theta^\star \geq 1} = \eval_\theta^\star \, 1_{\rho(\theta) \geq 1}. \]
Therefore, from the pointwise analysis above, 
\begin{align*}
m\{ \rho(\Theta), \eval_\Theta^\star(Z) \} & \overset{\text{\tiny (+)}}{\geq} \rho(\Theta) \, 1_{\rho(\Theta) < 1, \eval_\Theta^\star(Z) \geq 1} + \eval_\Theta^\star(Z) \, \rho(\Theta) \, 1_{\rho(\Theta) \geq 1, \eval_\Theta^\star(Z) \geq 1} \\
& = \rho(\Theta) \, 1_{\rho(\Theta) < 1} + \eval_\Theta^\star(Z) \, \rho(\Theta) \, 1_{\rho(\Theta) \geq 1}. 
\end{align*}
The (+) symbol above is to remind the reader that, while ``$\geq$'' holds pointwise, there's actually more that can be said.  That is, in addition to ``$\geq$'' for every $(z,\theta)$ pair, there exists a joint distribution for $(Z,\Theta)$, compatible with the available prior information, such that strict inequality holds with positive probability.  This implies that the inequality ``$\geq$'' highlighted with (+) is {\em strict inequality} ``$>$'' in expectation with respect to the aforementioned joint distribution.  Then 
\begin{align}
\uuprob\bigl[ m\{ \rho(\Theta), \eval_\Theta^\star(Z) \} \bigr] & > \uuprob\bigl\{ \rho(\Theta) \, 1_{\rho(\Theta) < 1} + \eval_\Theta^\star(Z) \, \rho(\Theta) \, 1_{\rho(\Theta) \geq 1} \bigr\} \notag \\
& = \sup_{\prior \in \credal} \E^{(Z,\Theta) \sim \prob_\bullet \otimes \prior} \bigl\{ \rho(\Theta) \, 1_{\rho(\Theta) < 1} + \eval_\Theta^\star(Z) \, \rho(\Theta) \, 1_{\rho(\Theta) \geq 1} \bigr\} \notag \\
& = \uprior\bigl\{ \rho(\Theta) \, 1_{\rho(\Theta) < 1} + \rho(\Theta) \, 1_{\rho(\Theta) \geq 1} \bigr\} \label{eq:dominance.a} \\
& = \uprior\,\rho \notag \\
& = 1, \label{eq:dominance.b}
\end{align}
where \eqref{eq:dominance.a} holds because $\eval^\star$ is an e-process relative to the model, i.e., $\E_\theta(\eval_\theta^\star) = 1$ for all $\theta$, and \eqref{eq:dominance.b} holds by the assumed admissibility of the regularizer $\rho$.  Therefore, the defining property \eqref{eq:merging} of an re-merging fails for the chosen $m$; that is, I constructed an e-process $(\eval^\star)$ such that $\uuprob[ m\{ \rho(\Theta), \eval_\Theta^\star(Z) \} ] > 1$.  Since $m$ was assumed to be re-merging, this creates the desired contradiction.  So, I conclude that, as claimed, there's no re-merging function $m$ that $\uuprob$-strictly weakly dominates the product rule in Equation~\eqref{eq:reg.eprocess}.  
\end{proof}

\end{document}
